% Generated draft: requires mathematical and source review.
\documentclass[11pt]{article}
\usepackage[margin=25mm]{geometry}
\usepackage{fontspec}
\setmainfont{Latin Modern Roman}
\usepackage{amsmath,amssymb,mathtools,unicode-math}
\setmathfont{Latin Modern Math}
\usepackage{enumitem,needspace,xurl,hyperref}
\setlength{\parindent}{0pt}
\setlength{\parskip}{5pt}
\setlength{\emergencystretch}{3em}
\newcommand{\N}{\mathbb{N}}
\newcommand{\Z}{\mathbb{Z}}
\newcommand{\Q}{\mathbb{Q}}
\newcommand{\R}{\mathbb{R}}
\newcommand{\C}{\mathbb{C}}
\newcommand{\F}{\mathbb{F}}
\newcommand{\A}{\mathbb{A}}
\newcommand{\PP}{\mathbb{P}}
\newcommand{\E}{\mathbb{E}}
\newcommand{\eps}{\varepsilon}
\newcommand{\dd}{\,\mathrm{d}}
\begin{document}
\section*{UnsolvedMath ID: 700443}
\textbf{Source problem code:} catalogue entry
\subsection{Definitions and mathematical statement}
Let \(p\) be a prime and let \(G\) be a finitely generated group. Let \(N\trianglelefteq G\) be a minimal normal subgroup: \(N\ne\{1\}\), and every subgroup \(M\trianglelefteq G\) with \(M\le N\) is either \(\{1\}\) or \(N\). Assume that \(N\) is an elementary abelian \(p\)-group, meaning that it is abelian and every nonidentity element has order \(p\); no finiteness assumption is imposed on \(N\).

For a finite generating set \(S\) of \(G\), define the ambient word length and the relative ball-counting function by
\[
\ell_S(g)=\min\{k\in\Z_{\ge0}:g=s_1\cdots s_k,\ s_i\in S\cup S^{-1}\},
\qquad
\gamma_{N,S}(n)=\bigl|\{x\in N:\ell_S(x)\le n\}\bigr|
\quad(n\in\Z_{\ge0}),
\]
where the empty product is the identity.

Is it true that, for every such \(p,G,N\) and every finite generating set \(S\) of \(G\), either \(N\) is finite or there exist constants \(C>0\), \(a>1\), and \(n_0\in\Z_{\ge0}\) such that
\[
\gamma_{N,S}(n)\ge C a^n
\qquad\text{for every integer }n\ge n_0?
\]
Here the source's unspecified ``growth function for the elements of \(N\)'' is interpreted as counting elements of \(N\) in ambient word balls, and ``bounded below by an exponential function'' is interpreted as the eventual lower bound displayed above. The constants may depend on \(G,N,S\). This exponential lower-bound property is independent of the choice of finite generating set, so the formulation also expresses the question for a chosen finite generating set.
\subsection{Short English statement}
For a finitely generated group \(G\), must every infinite minimal normal subgroup that is an elementary abelian \(p\)-group have exponentially many elements in the word balls of \(G\)? This uses ambient word length and an eventual exponential lower bound. The supplied registry labels the problem open; that status is not independently verified.
\subsection{Sources}
\begin{itemize}
\item Local UnsolvedMath registry, supplied complete source record: permanent ID \(700443\), exact source problem code \texttt{catalogue entry}, title ``Kourovka Notebook Problem 7.49.'' Supplied upstream revision: \texttt{d371ee917ac9647fbb06cd6fc529deec6f3dda82}.
\item The supplied attribution is to V. I. Trofimov, proposed in 1980, in \emph{The Kourovka Notebook (Unsolved Problems in Group Theory)}, 7th issue (1980), Problem 7.49. The record identifies the 21st issue edition (2026), \texttt{arXiv:1401.0300v48}, as its source edition.
\item The supplied collection description names E. I. Khukhro and V. D. Mazurov as editors of the Kourovka Notebook collection covering Issues 1--21 (1965--2026). The registry record was updated on 2026-10-16 and carries the status ``open''; this is supplied registry metadata, not an independently verified literature assessment.
\end{itemize}
\end{document}
